\section{Discussion, limitations, and conclusions}
\label{sec:discussion}
\subsection{Answers to the research questions}
\textbf{RQ1: representation and utility.} Low rank reduces shared payload and
can outperform the full private baseline in selected strong-privacy settings.
The four positive adjusted historical comparisons support that qualified
statement. Every nonprivate low-rank variant, however, underperforms the
converged full model, and low rank does not dominate across privacy targets.
The observed crossover combines capacity, optimization, clipping, and noise;
it is not an isolated dimensional effect.

\textbf{RQ2: fixed-factor effective noise.} FixedB supplies a useful linear
reference and smaller immediate perturbations in several diagnostics, but
obtains no positive adjusted primary comparison against Two. The simple
private popularity baseline exceeds all E1 collaborative point means at
targets at most four. Smaller coordinate or effective-matrix noise alone is
therefore an inadequate method-selection criterion in this benchmark.

\textbf{RQ3: persistent local state.} A perturbation to shared parameters can
leave local user vectors changed even after the shared state is reset. The
Full and FixedB probes show modest nonzero effects, as expected for coupled
local/global dynamics. The evidence does not establish a large general
relevance loss caused by this pathway.

\textbf{RQ4: validity of trajectory attribution.} In rebalanced two-factor
models, raw-coordinate common random numbers can produce much larger paired
distances than an aligned coupling. The continuation-without-balancing
control and ML-1M replication support a dependence on factor-basis changes.
This finding changes the interpretation of E2: the raw trajectory distance
cannot be read as a representation-invariant measure of an individual
noise draw's effect. It does not identify a unique causal coupling or a
new private optimizer.

\textbf{RQ5: ranking tails.} The exact conditional margin law can differ
substantially from a Gaussian in constructed regimes, but approximation
error is negligible in the preselected real checkpoint screen. The
non-Gaussian-tail candidate fails its practical advancement gate. The
result narrows the explanation of these experiments without claiming
Gaussianity throughout training.

\subsection{What the combined evidence explains}
The project demonstrates why an apparently intuitive story about private
recommendation can fail. Reducing shared dimension lowers a coordinate-noise
norm and payload, but also changes the learning geometry. Fixing a factor
removes a bilinear term, yet that term's measured energy can be small while
the method still loses utility. Large personalized score disturbances can
come from user scale, and large paired trajectory distances can arise from
coordinate coupling rather than a comparably large marginal ranking effect.

The relevant sequence of measurements is therefore: the bounded client
contribution, the effective shared signal and perturbation, the local user
state, score margins, candidate rankings, and relevance. Communication
should be measured both per exchange and over the full training schedule.
These measurements complement the accountant, which constrains information
release but does not predict usefulness. None is a substitute for a strong
simple comparator such as properly private popularity.

The most distinctive completed empirical observation is the sensitivity of
the personalized trajectory diagnostic to factor-basis coupling, supported
by the larger-data and balancing controls. Its components overlap existing
alignment and gauge methods. This report therefore treats it as a supported
diagnostic result whose paper-level novelty remains unconfirmed, not as an
established breakthrough or a new privacy theorem.

\subsection{Threats to validity}
\paragraph{Selection and benchmark exposure.}
ML-100K has repeated historical exposure, including the disclosed accidental
regression-test evaluations. Later freezes improve local experiment discipline
but do not undo accumulated development knowledge. E2b is an adaptive
follow-up, and E3/E4 questions were informed by earlier results. ML-1M
validation is now exposed; its test remains available for a future genuinely
frozen comparison.

\paragraph{Optimization and tuning.}
Historical B4 has more tuning than B2. E1 equalizes a finite candidate count,
not the attainable optimum of every method, and some selected rates are at
the boundary. Nonprivate convergence experiments and 50-round private runs
have different budgets. ML-1M transfer uses no retuning. These choices make
the recorded comparisons reproducible but do not prove method-independent
limits on privacy or rank.

\paragraph{Statistical interpretation.}
User-bootstrap intervals condition on retained training runs. Five seeds
provide limited information about rare training instabilities. E2/E3 probes
are correlated within checkpoints and are summarized descriptively. The
small observer subset in E3 limits resolution of validation differences.
Activity, cold-item, and personalization ablations are exploratory rather
than multiplicity-adjusted subgroup discoveries.

\paragraph{Privacy and systems scope.}
The simulator assumes only noised aggregate visibility and implements no
secure aggregation. Public population/catalog assumptions, nonprivate
selection, internal checkpoints, and diagnostic releases fall outside the
per-run guarantee. The modeled communication excludes protocol overhead
and client availability costs. It does not evaluate a deployed mobile system.

\paragraph{Generalization and numerical scope.}
Both benchmarks are MovieLens variants, not independent services. Per-user
chronological splitting is not a global future-data prediction experiment.
There are very few cold targets. Effective-noise formulas condition on a
single state, while E4 samples a finite collection of round-50 margins.
Large finite local vectors and float32 update arithmetic require careful
numerical interpretation. The report makes no claim of universal bitwise
reproducibility across environments.

\subsection{Future work justified by the evidence}
The first priority is a faithful published private ALS or federated-MF
comparator. Such a comparison requires matching the privacy unit, adjacency,
trust model, feedback objective, candidate filter, and communication
definition. Perturbing sufficient statistics rather than BPR updates may
change the relevant signal-to-noise regime. The literature/implementation
audit is complete enough to specify this next comparison, but no such
baseline result is included as completed work.

The second priority is metadata-only local personalization, followed by a
small private collaborative residual. Public genres or other permitted
item features could supply useful structure before any private collaborative
update is spent. A meaningful extension would test incremental residual
value against that public-only control and private popularity under an
explicit total communication cap. Cold items and low-activity users are
natural targets, provided the dataset supplies enough examples for reliable
analysis. The underlying public-feature idea already has prior work.

A third possibility is to test whether coupling-aware diagnostics change
a consequential scientific conclusion outside the present pilot. The
target conclusion and decision criterion should be frozen before running
more probes. Merely repeating the same alignment contrast at additional
amplitudes would add volume without necessarily establishing novelty.
A useful method would also need to avoid counterfactual access to private
reference states and receive fresh privacy accounting.

\subsection{Conclusion}
This project began with low-dimensional sharing as a possible route to
better private federated recommendation and developed into a systematic
study of how parameter noise becomes prediction disturbance. The final
record includes a reproducible benchmark, controlled fixed-factor and
private-popularity comparisons, conditional noise analysis, pulse/reset
experiments, a larger-data replication, and a rejected ranking-tail
explanation. The evidence supports conditional low-rank gains and a
specific limitation of factor-coordinate trajectory measurements, while
rejecting broader claims that fixed factors, smaller coordinate norms,
or bilinear tails alone solve the observed utility problem.

The principal outcome is a more precise experimental foundation for the
next research decision. It preserves both favorable and unfavorable
results, defines the scope of privacy and uncertainty, and identifies
stronger comparators and public-information controls needed before claiming
a new recommendation method. That foundation, rather than an unsupported
novelty claim, is the completed contribution of this undergraduate project.
